% ══════════════════════════════════════════════════════════════════════════════
%                         فصل اول
%                 زایش طیف سیاسی: از انقلاب فرانسه تا قرن نوزدهم
% ══════════════════════════════════════════════════════════════════════════════

\chapter{زایش طیف سیاسی}
\label{ch:birth-of-spectrum}

\begin{flushright}
\textit{«در آن روز، جایی که نشستی، سرنوشت تاریخ را رقم زد.»}
\end{flushright}

% ══════════════════════════════════════════════════════════════════════════════
%                         مقدمه فصل
% ══════════════════════════════════════════════════════════════════════════════

\lettrine[lines=3, lraise=0.1, nindent=0.5em]{\textcolor{LeftRed}{ا}}{مروزه} 
وقتی از «چپ» و «راست» سخن می‌گوییم، گویی از مفاهیمی ازلی و ابدی سخن می‌گوییم. اما این واژگان عمری کوتاه‌تر از دو قرن و نیم دارند. پیش از تابستان ۱۷۸۹، هیچ‌کس در هیچ زبانی از «سیاست چپ» یا «جناح راست» سخن نمی‌گفت. این فصل داستان زایش یکی از مهم‌ترین استعاره‌های سیاسی بشر را روایت می‌کند.

% ══════════════════════════════════════════════════════════════════════════════
\section{پیش از چپ و راست: جهان سیاسی پیشامدرن}
% ══════════════════════════════════════════════════════════════════════════════

پیش از عصر انقلاب‌ها، جهان سیاسی به شکلی متفاوت سازمان‌دهی می‌شد. اختلافات سیاسی نه بر محور چپ و راست، بلکه بر محورهای دیگری شکل می‌گرفت:

\subsection{نظام‌های سلسله‌مراتبی}

در اروپای قرون وسطی و اوایل دوران مدرن، جامعه بر اساس «طبقات» یا «صنوف» (\term{Estates}) سازمان‌دهی می‌شد:

\begin{itemize}[nosep, rightmargin=1.5cm]
    \item \textbf{طبقه اول:} روحانیون (کلیسا)
    \item \textbf{طبقه دوم:} اشراف (نجبا)
    \item \textbf{طبقه سوم:} عوام (بورژوازی، پیشه‌وران، دهقانان)
\end{itemize}

\begin{historybox}[نظام سه‌گانه صنوف]
در فرانسه پیش از انقلاب، جمعیت به سه صنف تقسیم می‌شد. روحانیون (حدود ۱۰۰ هزار نفر) و اشراف (حدود ۴۰۰ هزار نفر) از امتیازات ویژه برخوردار بودند و مالیات نمی‌پرداختند، در حالی که طبقه سوم (حدود ۲۶ میلیون نفر) بار اصلی مالیات‌ها را بر دوش می‌کشید.
\end{historybox}

\subsection{غیاب مفهوم طیف سیاسی}

در این نظام، اختلافات سیاسی عمدتاً حول محورهای زیر بود:

\begin{enumerate}[nosep, rightmargin=1.5cm]
    \item \textbf{وفاداری به شخص:} حمایت از این یا آن شاهزاده، این یا آن خاندان
    \item \textbf{اختلافات مذهبی:} کاتولیک در برابر پروتستان
    \item \textbf{منافع صنفی:} منافع اشراف در برابر منافع شهرنشینان
    \item \textbf{محلی‌گرایی:} منافع این منطقه در برابر آن منطقه
\end{enumerate}

\begin{defbox}[تعریف: پیشامدرنیته سیاسی]
دوران پیشامدرن سیاسی به دورانی اطلاق می‌شود که در آن:
\begin{itemize}[nosep]
    \item مشروعیت قدرت «از بالا» (خدا، سنت) می‌آمد، نه «از پایین» (مردم)
    \item مفهوم «شهروند» وجود نداشت؛ افراد «رعیت» بودند
    \item ایده‌ی تغییر بنیادین نظام سیاسی نامتصور بود
\end{itemize}
\end{defbox}

% ══════════════════════════════════════════════════════════════════════════════
\section{انقلاب فرانسه: لحظه زایش}
% ══════════════════════════════════════════════════════════════════════════════

\subsection{بحران مالی و تشکیل مجلس عمومی}

در سال ۱۷۸۹، فرانسه در بحران مالی عمیقی دست‌وپا می‌زد. جنگ‌های پرهزینه، از جمله حمایت از انقلاب آمریکا، خزانه را خالی کرده بود. \person{لویی شانزدهم}، پادشاه فرانسه، چاره‌ای جز فراخواندن \term{مجلس عمومی} (\term{États généraux}) ندید—نهادی که از سال ۱۶۱۴ تشکیل نشده بود.

\begin{figure}[htbp]
\centering
\begin{tikzpicture}[
    node distance=1.5cm,
    every node/.style={font=\small}
]

% عنوان
\node[font=\large\bfseries, text=DarkGray] at (0, 5) {\fa{تایم‌لاین: از بحران تا انقلاب}};

% خط زمان
\draw[line width=3pt, MediumGray] (-7, 3.5) -- (7, 3.5);

% رویدادها
% ۱۷۸۸
\fill[RightBlue] (-6, 3.5) circle (8pt);
\node[below, text width=2.5cm, align=center, text=DarkGray] at (-6, 2.8) {
    \fa{اوت ۱۷۸۸}\\
    \fa{فراخوان}\\
    \fa{مجلس عمومی}
};

% ۵ مه ۱۷۸۹
\fill[MediumGray] (-3.5, 3.5) circle (8pt);
\node[below, text width=2.5cm, align=center, text=DarkGray] at (-3.5, 2.8) {
    \fa{۵ مه ۱۷۸۹}\\
    \fa{افتتاح}\\
    \fa{مجلس عمومی}
};

% ۱۷ ژوئن
\fill[CenterGreen] (-1, 3.5) circle (8pt);
\node[below, text width=2.5cm, align=center, text=DarkGray] at (-1, 2.8) {
    \fa{۱۷ ژوئن}\\
    \fa{اعلام}\\
    \fa{مجلس ملی}
};

% ۲۰ ژوئن
\fill[GoldAccent] (1.5, 3.5) circle (8pt);
\node[below, text width=2.5cm, align=center, text=DarkGray] at (1.5, 2.8) {
    \fa{۲۰ ژوئن}\\
    \fa{سوگند}\\
    \fa{زمین تنیس}
};

% ۱۴ ژوئیه
\fill[LeftRed] (4, 3.5) circle (8pt);
\node[below, text width=2.5cm, align=center, text=DarkGray] at (4, 2.8) {
    \fa{۱۴ ژوئیه}\\
    \fa{سقوط}\\
    \fa{باستیل}
};

% ۲۶ اوت
\fill[PurpleAccent] (6.5, 3.5) circle (8pt);
\node[below, text width=2.5cm, align=center, text=DarkGray] at (6.5, 2.8) {
    \fa{۲۶ اوت}\\
    \fa{اعلامیه}\\
    \fa{حقوق بشر}
};

\end{tikzpicture}
\caption{رویدادهای کلیدی آغاز انقلاب فرانسه (۱۷۸۸-۱۷۸۹)}
\label{fig:french-rev-timeline}
\end{figure}

\subsection{از مجلس عمومی به مجلس ملی}

مجلس عمومی در ۵ مه ۱۷۸۹ در ورسای گشایش یافت. اما خیلی زود اختلاف بر سر شیوه رأی‌گیری بالا گرفت:

\begin{debatebox}[رأی‌گیری: صنفی یا فردی؟]
\textbf{اشراف و روحانیون:} هر صنف یک رأی داشته باشد (۲ رأی در برابر ۱)

\textbf{طبقه سوم:} رأی‌گیری سرشماری (هر نماینده یک رأی)

این اختلاف ظاهراً فنی، در واقع اختلافی بنیادین بود: آیا قدرت متعلق به صنوف ممتاز است یا به «ملت»؟
\end{debatebox}

در ۱۷ ژوئن، نمایندگان طبقه سوم، به همراه برخی روحانیون و اشراف لیبرال، خود را \term{مجلس ملی} خواندند. این لحظه، لحظه‌ای انقلابی بود: برای نخستین بار، «ملت» به عنوان منبع مشروعیت معرفی شد.

\subsection{تولد چپ و راست: سپتامبر ۱۷۸۹}

پس از سقوط باستیل و رویدادهای پرشتاب تابستان، مجلس ملی مؤسسان مشغول بحث درباره ساختار قانون اساسی جدید بود. یکی از مسائل کلیدی، \textbf{حق وتوی شاه} بود.

\begin{historybox}[لحظه تاریخی: چینش مجلس]
در جلسات سپتامبر ۱۷۸۹، نمایندگانی که از حق وتوی مطلق شاه حمایت می‌کردند، در \textbf{سمت راست} رئیس مجلس نشستند. مخالفان این حق در \textbf{سمت چپ} قرار گرفتند. این چینش تصادفی، به استعاره‌ای ماندگار تبدیل شد.
\end{historybox}

\begin{figure}[htbp]
\centering
\begin{tikzpicture}[scale=0.9]

% عنوان
\node[font=\large\bfseries, text=DarkGray] at (0, 7) {\fa{چینش صندلی‌ها در مجلس ملی مؤسسان (۱۷۸۹)}};

% نیم‌دایره مجلس
\draw[line width=2pt, color=DarkGray, fill=VeryLightGray] 
    (0, 0) -- (-7, 0) arc (180:0:7) -- cycle;

% بخش‌بندی
% راست (سلطنت‌طلبان)
\fill[RightBlue, opacity=0.6] 
    (0, 0) -- (7, 0) arc (0:60:7) -- cycle;
\node[text=white, font=\bfseries] at (4.5, 2.5) {
    \begin{tabular}{c}
    \fa{راست}\\
    \fa{سلطنت‌طلبان}
    \end{tabular}
};

% راست میانه (مشروطه‌خواهان محافظه‌کار)
\fill[RightBlueLight, opacity=0.6] 
    (0, 0) -- (60:7) arc (60:90:7) -- cycle;
\node[text=DarkGray, font=\small\bfseries] at (2, 5.5) {
    \begin{tabular}{c}
    \fa{مشروطه‌خواهان}\\
    \fa{محافظه‌کار}
    \end{tabular}
};

% میانه
\fill[CenterGreen, opacity=0.5] 
    (0, 0) -- (90:7) arc (90:110:7) -- cycle;
\node[text=white, font=\small\bfseries, rotate=-10] at (-0.8, 6) {\fa{میانه}};

% چپ میانه (مشروطه‌خواهان لیبرال)
\fill[CenterGreenLight, opacity=0.6] 
    (0, 0) -- (110:7) arc (110:140:7) -- cycle;
\node[text=DarkGray, font=\small\bfseries] at (-3.5, 5) {
    \begin{tabular}{c}
    \fa{لیبرال‌ها}
    \end{tabular}
};

% چپ (جمهوری‌خواهان و رادیکال‌ها)
\fill[LeftRed, opacity=0.6] 
    (0, 0) -- (140:7) arc (140:180:7) -- cycle;
\node[text=white, font=\bfseries] at (-5.5, 2.5) {
    \begin{tabular}{c}
    \fa{چپ}\\
    \fa{رادیکال‌ها}
    \end{tabular}
};

% جایگاه رئیس
\fill[GoldAccent] (0, -0.5) ellipse (1.2 and 0.5);
\node[text=white, font=\small\bfseries] at (0, -0.5) {\fa{رئیس مجلس}};

% فلش‌های راهنما
\draw[thickarrow, color=RightBlue] (8, 3) -- (6, 3);
\node[left, text=RightBlue, font=\small] at (9.5, 3) {\fa{هوادار شاه}};

\draw[thickarrow, color=LeftRed] (-8, 3) -- (-6, 3);
\node[right, text=LeftRed, font=\small] at (-9.5, 3) {\fa{مخالف شاه}};

% راهنمای پایین
\node[text=MediumGray, font=\footnotesize] at (0, -2) {
    \fa{این چینش فیزیکی، مبنای استعاره «چپ و راست» در سیاست شد.}
};

\end{tikzpicture}
\caption{نمودار چینش مجلس ملی مؤسسان فرانسه}
\label{fig:assembly-seating}
\end{figure}

\subsection{محتوای اولیه چپ و راست}

در آن روزهای نخستین، تمایز چپ و راست معنایی نسبتاً ساده داشت:

\begin{table}[htbp]
\centering
\caption{معنای اولیه چپ و راست (۱۷۸۹-۱۷۹۱)}
\label{tab:initial-left-right}
\begin{tabular}{>{\columncolor{RightBlueLight!30}}r | >{\columncolor{LeftRedLight!30}}r}
\toprule
\textbf{راست} & \textbf{چپ} \\
\midrule
حفظ اختیارات شاه & محدودسازی شاه \\
سلطنت مطلقه یا نیمه‌مطلقه & سلطنت مشروطه یا جمهوری \\
حفظ امتیازات کلیسا & جدایی دین از دولت \\
تغییرات تدریجی & تغییرات بنیادین \\
نظم و سنت & آزادی و برابری \\
\bottomrule
\end{tabular}
\end{table}

% ══════════════════════════════════════════════════════════════════════════════
\section{رادیکالی‌شدن انقلاب}
% ══════════════════════════════════════════════════════════════════════════════

انقلاب فرانسه در مسیر خود به شدت رادیکال شد. این رادیکالی‌شدن، طیف سیاسی را نیز گسترش داد.

\subsection{از ژیروندن تا ژاکوبن}

\begin{figure}[htbp]
\centering
\begin{tikzpicture}[scale=0.85]

% عنوان
\node[font=\large\bfseries, text=DarkGray] at (0, 6) {\fa{تایم‌لاین رادیکالی‌شدن انقلاب (۱۷۸۹-۱۷۹۴)}};

% خط زمان اصلی
\draw[line width=4pt, MediumGray] (-7, 4) -- (7, 4);

% مراحل
% مرحله ۱: لیبرال
\fill[CenterGreen] (-5.5, 4) circle (10pt);
\node[above, text=CenterGreen, font=\bfseries] at (-5.5, 4.5) {\fa{۱۷۸۹-۹۱}};
\node[below, text width=2.5cm, align=center, text=DarkGray, font=\small] at (-5.5, 3.2) {
    \fa{مرحله لیبرال}\\
    \fa{مشروطه‌خواهی}
};

% مرحله ۲: ژیروندن
\fill[OrangeAccent] (-2, 4) circle (10pt);
\node[above, text=OrangeAccent, font=\bfseries] at (-2, 4.5) {\fa{۱۷۹۱-۹۲}};
\node[below, text width=2.5cm, align=center, text=DarkGray, font=\small] at (-2, 3.2) {
    \fa{ژیروندن‌ها}\\
    \fa{جمهوری میانه‌رو}
};

% مرحله ۳: ژاکوبن
\fill[LeftRed] (1.5, 4) circle (10pt);
\node[above, text=LeftRed, font=\bfseries] at (1.5, 4.5) {\fa{۱۷۹۲-۹۴}};
\node[below, text width=2.5cm, align=center, text=DarkGray, font=\small] at (1.5, 3.2) {
    \fa{ژاکوبن‌ها}\\
    \fa{دوران ترور}
};

% مرحله ۴: ترمیدور
\fill[RightBlueLight] (5, 4) circle (10pt);
\node[above, text=RightBlueLight, font=\bfseries] at (5, 4.5) {\fa{۱۷۹۴-۹۹}};
\node[below, text width=2.5cm, align=center, text=DarkGray, font=\small] at (5, 3.2) {
    \fa{ترمیدور}\\
    \fa{بازگشت به میانه}
};

% فلش نشان‌دهنده رادیکالی‌شدن
\draw[thickarrow, color=LeftRed, opacity=0.7] (-6, 2) -- (2, 2);
\node[text=LeftRed, font=\small] at (-2, 1.5) {\fa{رادیکالی‌شدن}};

% فلش واکنش
\draw[thickarrow, color=RightBlue, opacity=0.7] (2.5, 2) -- (5.5, 2);
\node[text=RightBlue, font=\small] at (4, 1.5) {\fa{واکنش}};

\end{tikzpicture}
\caption{سیر رادیکالی‌شدن و واکنش در انقلاب فرانسه}
\label{fig:radicalization}
\end{figure}

\begin{defbox}[تعریف: ژیروندن و ژاکوبن]
\textbf{ژیروندن‌ها} (\term{Girondins}): جناح میانه‌رو جمهوری‌خواه، عمدتاً از بورژوازی شهرستانی، طرفدار جمهوری فدرالی و جنگ با اروپا.

\textbf{ژاکوبن‌ها} (\term{Jacobins}): جناح رادیکال، طرفدار جمهوری متمرکز، اصلاحات اجتماعی عمیق، و سرکوب مخالفان. نام‌شان از دیرِ ژاکوبن، محل گردهمایی‌شان، گرفته شده.
\end{defbox}

\subsection{دوران ترور و پیامدهایش}

دوران ترور (۱۷۹۳-۱۷۹۴) تحت رهبری \person{روبسپیر} و کمیته نجات ملی، نشان داد که انقلاب می‌تواند فرزندان خود را نیز ببلعد. هزاران نفر، از جمله بسیاری از انقلابیون اولیه، گیوتین شدند.

\begin{quotebox}[ورنیو، رهبر ژیروندن، پیش از اعدام]
\textit{«انقلاب همچون ساتورن، فرزندان خود را می‌بلعد.»}
\end{quotebox}

این تجربه تلخ، برای همیشه در حافظه سیاسی اروپا حک شد و یکی از مهم‌ترین انتقادات به رادیکالیسم چپ شد: ترس از اینکه انقلاب به استبداد جدید بینجامد.

% ══════════════════════════════════════════════════════════════════════════════
\section{شخصیت‌های کلیدی}
% ══════════════════════════════════════════════════════════════════════════════

درک زایش طیف سیاسی، بدون آشنایی با شخصیت‌های کلیدی ممکن نیست. در اینجا دو چهره‌ی نمادین معرفی می‌شوند که به نوعی «پدران» راست و چپ مدرن به شمار می‌آیند.

\begin{personbox}[ادموند برک (۱۷۲۹-۱۷۹۷)]

\textbf{ملیت:} ایرلندی-بریتانیایی

\textbf{اهمیت:} پدر محافظه‌کاری مدرن

\textbf{اثر کلیدی:} \book{تأملاتی درباره انقلاب فرانسه} (۱۷۹۰)

\mediumspace

\person{ادموند برک}، نماینده پارلمان بریتانیا، در واکنش به انقلاب فرانسه کتابی نوشت که بنیان‌گذار اندیشه محافظه‌کار شد. او استدلال کرد که:

\begin{itemize}[nosep, rightmargin=0.5cm]
    \item \textbf{خرد جمعی سنت} برتر از خرد فردی فیلسوفان است
    \item نهادهای اجتماعی محصول قرن‌ها آزمون‌وخطا هستند
    \item تغییرات باید \textbf{تدریجی} و محافظه‌کارانه باشند
    \item انقلاب‌های رادیکال به هرج‌ومرج و استبداد می‌انجامند
\end{itemize}

\mediumspace

\textbf{نقل‌قول معروف:} \textit{«جامعه قراردادی است میان مردگان، زندگان و آنان که هنوز زاده نشده‌اند.»}
\end{personbox}

\vspace{15pt}

\begin{personbox}[توماس پین (۱۷۳۷-۱۸۰۹)]

\textbf{ملیت:} بریتانیایی-آمریکایی

\textbf{اهمیت:} نظریه‌پرداز رادیکالیسم دموکراتیک

\textbf{آثار کلیدی:} \book{عقل سلیم} (۱۷۷۶)، \book{حقوق بشر} (۱۷۹۱)

\mediumspace

\person{توماس پین} در پاسخ به برک، کتاب \book{حقوق بشر} را نوشت. او استدلال کرد که:

\begin{itemize}[nosep, rightmargin=0.5cm]
    \item هر نسل \textbf{حق تعیین سرنوشت} خود را دارد
    \item سنت نمی‌تواند توجیه‌گر بی‌عدالتی باشد
    \item حکومت باید بر رضایت مردم استوار باشد
    \item \textbf{حقوق طبیعی} (آزادی، برابری) مقدم بر هر قانونی است
\end{itemize}

\mediumspace

\textbf{نقل‌قول معروف:} \textit{«زمانی که همه دروغ باور شوند، گفتن حقیقت عملی انقلابی است.»}
\end{personbox}

\begin{comparebox}[مقایسه: برک در برابر پین]
\begin{table}[H]
\centering
\begin{tabular}{r|c|c}
\toprule
\textbf{موضوع} & \textbf{برک (راست)} & \textbf{پین (چپ)} \\
\midrule
منبع مشروعیت & سنت و تجربه تاریخی & حقوق طبیعی و عقل \\
نگاه به انسان & ناقص و نیازمند نهادها & قابل تکامل و خودمختار \\
تغییر & تدریجی و محتاطانه & بنیادین اگر لازم باشد \\
انقلاب فرانسه & فاجعه‌بار و خطرناک & ضروری و الهام‌بخش \\
سلطنت & مشروطه ضروری است & جمهوری برتر است \\
\bottomrule
\end{tabular}
\end{table}
\end{comparebox}

% ══════════════════════════════════════════════════════════════════════════════
\section{گسترش مفهوم در قرن نوزدهم}
% ══════════════════════════════════════════════════════════════════════════════

پس از انقلاب فرانسه و دوران ناپلئون، مفهوم چپ و راست به تدریج در سراسر اروپا گسترش یافت و محتوای آن نیز تحول پیدا کرد.

\subsection{بازگشت و انقلاب: ۱۸۱۵-۱۸۴۸}

\begin{historybox}[کنگره وین (۱۸۱۴-۱۸۱۵)]
پس از شکست ناپلئون، قدرت‌های اروپایی در کنگره وین گرد آمدند تا نظم پیش از انقلاب را احیا کنند. این دوره «بازگشت» (\term{Restoration}) نامیده می‌شود. اما ایده‌های انقلاب را نمی‌شد بازگرداند.
\end{historybox}

در این دوره، طیف سیاسی شفاف‌تر شد:

\begin{itemize}[nosep, rightmargin=1.5cm]
    \item \textbf{راست افراطی (ارتجاعی):} بازگشت کامل به نظم پیشین، سلطنت مطلقه
    \item \textbf{راست میانه (محافظه‌کار):} سلطنت مشروطه محدود
    \item \textbf{میانه (لیبرال):} قانون اساسی، حقوق فردی، بازار آزاد
    \item \textbf{چپ میانه (رادیکال):} جمهوری، حق رأی گسترده
    \item \textbf{چپ افراطی (سوسیالیست):} دموکراسی اجتماعی، اصلاحات اقتصادی
\end{itemize}

\subsection{انقلاب‌های ۱۸۴۸: بهار ملت‌ها}

سال ۱۸۴۸ شاهد موج انقلابی در سراسر اروپا بود. این انقلاب‌ها، هرچند اغلب شکست خوردند، طیف سیاسی را به شکل جدیدی ترسیم کردند.

\begin{figure}[htbp]
\centering
\begin{tikzpicture}[scale=0.8]

% عنوان
\node[font=\large\bfseries, text=DarkGray] at (0, 7.5) {\fa{انقلاب‌های ۱۸۴۸ در اروپا}};

% خط زمان
\draw[line width=3pt, MediumGray] (-7, 5.5) -- (7, 5.5);

\node[above, text=MediumGray, font=\bfseries] at (-7, 5.8) {\fa{ژانویه}};
\node[above, text=MediumGray, font=\bfseries] at (7, 5.8) {\fa{دسامبر}};

% رویدادها
% سیسیل
\fill[OrangeAccent] (-6, 5.5) circle (8pt);
\node[below, text width=2cm, align=center, font=\small] at (-6, 4.8) {
    \fa{ژانویه}\\
    \fa{سیسیل}
};

% فرانسه
\fill[LeftRed] (-4, 5.5) circle (10pt);
\node[below, text width=2cm, align=center, font=\small] at (-4, 4.8) {
    \fa{فوریه}\\
    \fa{پاریس}
};

% آلمان/اتریش
\fill[GoldAccent] (-1.5, 5.5) circle (9pt);
\node[below, text width=2cm, align=center, font=\small] at (-1.5, 4.8) {
    \fa{مارس}\\
    \fa{وین، برلین}
};

% مجارستان
\fill[CenterGreen] (1, 5.5) circle (8pt);
\node[below, text width=2cm, align=center, font=\small] at (1, 4.8) {
    \fa{مارس}\\
    \fa{بوداپست}
};

% ایتالیا
\fill[OrangeAccent] (3.5, 5.5) circle (8pt);
\node[below, text width=2cm, align=center, font=\small] at (3.5, 4.8) {
    \fa{مارس-آوریل}\\
    \fa{میلان، ونیز}
};

% سرکوب
\fill[RightBlueDark] (6, 5.5) circle (8pt);
\node[below, text width=2cm, align=center, font=\small] at (6, 4.8) {
    \fa{پاییز}\\
    \fa{سرکوب}
};

% خط اهمیت
\node[text=LeftRed, font=\bfseries] at (-4, 3.5) {\fa{مانیفست کمونیست منتشر شد}};
\draw[dashedarrow, color=LeftRed] (-4, 3.8) -- (-4, 5);

\end{tikzpicture}
\caption{موج انقلابی ۱۸۴۸ در اروپا}
\label{fig:1848-revolutions}
\end{figure}

\subsection{ظهور سوسیالیسم: بُعد جدید}

مهم‌ترین تحول قرن نوزدهم، ظهور \keyword{سوسیالیسم} بود. این جنبش، پرسش‌های جدیدی را به طیف سیاسی افزود:

\begin{defbox}[تعریف: سوسیالیسم (معنای اولیه)]
در قرن نوزدهم، سوسیالیسم به مجموعه‌ای از ایده‌ها و جنبش‌ها اطلاق می‌شد که خواستار:
\begin{itemize}[nosep]
    \item مالکیت جمعی یا اجتماعی ابزار تولید
    \item پایان استثمار کار توسط سرمایه
    \item برابری اقتصادی و اجتماعی
    \item حمایت از طبقه کارگر
\end{itemize}
بودند. سوسیالیسم در اشکال مختلفی (تخیلی، علمی، آنارشیستی) ظاهر شد.
\end{defbox}

با ظهور سوسیالیسم، طیف سیاسی از یک پرسش صرفاً سیاسی (چه کسی حکومت کند؟) به پرسشی اقتصادی-اجتماعی (مالکیت و توزیع ثروت چگونه باشد؟) گسترش یافت.

% ══════════════════════════════════════════════════════════════════════════════
\section{اختلافات اولیه درون هر جناح}
% ══════════════════════════════════════════════════════════════════════════════

از همان آغاز، نه چپ یکدست بود و نه راست. درک این تنوع درونی برای فهم تاریخ سیاسی ضروری است.

\subsection{شکاف‌های درون راست}

\begin{table}[htbp]
\centering
\caption{گرایش‌های مختلف در «راست» قرن نوزدهم}
\label{tab:right-factions}
\begin{tabular}{r|r|r}
\toprule
\textbf{گرایش} & \textbf{خواسته اصلی} & \textbf{نمونه} \\
\midrule
ارتجاعی & بازگشت به نظم پیشین & اولترارویالیست‌ها \\
محافظه‌کار سنتی & حفظ سلطنت مشروطه & توری‌های بریتانیا \\
لیبرال محافظه‌کار & بازار آزاد + نظم اجتماعی & ویگ‌های بریتانیا \\
ناسیونالیست & وحدت ملی و قدرت & بیسمارک \\
\bottomrule
\end{tabular}
\end{table}

\begin{debatebox}[اختلاف کلیدی راست: سنت یا بازار؟]
آیا «راست» باید از سنت‌ها و نهادهای تاریخی (کلیسا، اشرافیت، جامعه روستایی) دفاع کند، حتی اگر با منطق بازار آزاد تعارض داشته باشند؟

\textbf{محافظه‌کاران سنتی:} بله، بازار ابزار است نه هدف.

\textbf{لیبرال‌های اقتصادی:} خیر، آزادی اقتصادی مقدم است.

این تنش تا امروز در درون راست باقی است.
\end{debatebox}

\subsection{شکاف‌های درون چپ}

\begin{table}[htbp]
\centering
\caption{گرایش‌های مختلف در «چپ» قرن نوزدهم}
\label{tab:left-factions}
\begin{tabular}{r|r|r}
\toprule
\textbf{گرایش} & \textbf{خواسته اصلی} & \textbf{نمونه} \\
\midrule
لیبرال رادیکال & جمهوری و حقوق مدنی & رادیکال‌های فرانسوی \\
دموکرات & حق رأی عمومی & چارتیست‌های انگلیس \\
سوسیالیست تخیلی & جوامع نمونه & سن‌سیمون، فوریه \\
سوسیالیست علمی & انقلاب پرولتری & مارکس، انگلس \\
آنارشیست & الغای دولت & پرودون، باکونین \\
\bottomrule
\end{tabular}
\end{table}

\begin{debatebox}[اختلاف کلیدی چپ: دولت یا ضد دولت؟]
آیا چپ باید از دولت قوی برای تحقق برابری استفاده کند، یا دولت خود ابزار سرکوب است؟

\textbf{سوسیالیست‌های دولت‌گرا:} دولت کارگری ابزار رهایی است.

\textbf{آنارشیست‌ها:} هر دولتی، حتی کارگری، به سرکوب می‌انجامد.

این اختلاف در انترناسیونال اول به انشعاب انجامید.
\end{debatebox}

% ══════════════════════════════════════════════════════════════════════════════
\section{نقد مفهوم طیف خطی}
% ══════════════════════════════════════════════════════════════════════════════

از همان آغاز، برخی متفکران به محدودیت‌های مدل خطی چپ-راست اشاره کردند.

\subsection{محدودیت‌های مدل یک‌بعدی}

\begin{enumerate}[nosep, rightmargin=1.5cm]
    \item \textbf{ساده‌سازی بیش از حد:} مسائل پیچیده در یک خط جا نمی‌گیرند
    \item \textbf{ناهمگونی درونی:} چپ‌ها و راست‌های مختلف ممکن است در مسائلی همراه باشند
    \item \textbf{مسائل متعامد:} آزادی اقتصادی و آزادی فردی لزوماً همراه نیستند
    \item \textbf{تاریخ‌مندی:} معنای چپ و راست در زمان‌ها و مکان‌های مختلف متفاوت است
\end{enumerate}

\begin{figure}[htbp]
\centering
\begin{tikzpicture}[scale=0.85]

% عنوان
\node[font=\large\bfseries, text=DarkGray] at (0, 6) {\fa{مدل خطی و محدودیت‌هایش}};

% مدل خطی ساده
\draw[line width=3pt, MediumGray] (-6, 4) -- (6, 4);
\fill[LeftRed] (-5, 4) circle (10pt);
\node[below, text=LeftRed, font=\bfseries] at (-5, 3.4) {\fa{چپ افراطی}};

\fill[LeftRedLight] (-2.5, 4) circle (8pt);
\node[below, text=LeftRedLight, font=\small] at (-2.5, 3.5) {\fa{چپ}};

\fill[CenterGreen] (0, 4) circle (8pt);
\node[below, text=CenterGreen, font=\small] at (0, 3.5) {\fa{میانه}};

\fill[RightBlueLight] (2.5, 4) circle (8pt);
\node[below, text=RightBlueLight, font=\small] at (2.5, 3.5) {\fa{راست}};

\fill[RightBlue] (5, 4) circle (10pt);
\node[below, text=RightBlue, font=\bfseries] at (5, 3.4) {\fa{راست افراطی}};

% علامت سؤال
\node[text=OrangeAccent, font=\Huge\bfseries] at (0, 1.5) {?};

% سؤالات
\node[text=DarkGray, font=\small, text width=10cm, align=center] at (0, 0.3) {
    \fa{لیبرتارین‌ها کجا قرار می‌گیرند؟ (آزادی اقتصادی + آزادی فردی)}
};
\node[text=DarkGray, font=\small, text width=10cm, align=center] at (0, -0.5) {
    \fa{فاشیست‌ها چپ‌اند یا راست؟ (ضد سرمایه‌داری + ضد سوسیالیسم)}
};
\node[text=DarkGray, font=\small, text width=10cm, align=center] at (0, -1.3) {
    \fa{محیط‌زیست‌گرایان کجایند؟ (نه کاملاً چپ، نه کاملاً راست)}
};

\end{tikzpicture}
\caption{محدودیت‌های مدل خطی طیف سیاسی}
\label{fig:linear-limitations}
\end{figure}

\subsection{پیش‌درآمد مدل‌های جایگزین}

این نقدها، بعدها به مدل‌های پیچیده‌تر انجامید که در فصل آخر کتاب بررسی خواهند شد:

\begin{itemize}[nosep, rightmargin=1.5cm]
    \item \textbf{مدل دوبعدی:} افزودن محور آزادی‌خواهی/اقتدارگرایی
    \item \textbf{مدل نعل اسب:} افراط‌های چپ و راست به هم نزدیک‌اند
    \item \textbf{مدل‌های چندمحوری:} هر مسئله محور جداگانه‌ای دارد
\end{itemize}

با این حال، مدل ساده چپ-راست همچنان پرکاربردترین ابزار برای فهم سیاست است—شاید دقیقاً به دلیل سادگی‌اش.

% ══════════════════════════════════════════════════════════════════════════════
%                         خلاصه و پرسش‌ها
% ══════════════════════════════════════════════════════════════════════════════

\newpage

\begin{summarybox}

\textbf{۱. زایش مفهوم:} چپ و راست در انقلاب فرانسه ۱۷۸۹ زاده شدند، از چینش فیزیکی نمایندگان در مجلس.

\textbf{۲. معنای اولیه:} راست = حفظ اختیارات شاه و نظم موجود؛ چپ = محدودسازی قدرت و تغییر.

\textbf{۳. رادیکالی‌شدن:} انقلاب فرانسه از مشروطه‌خواهی لیبرال به جمهوری‌خواهی رادیکال و دوران ترور رسید.

\textbf{۴. شخصیت‌های کلیدی:} برک (پدر محافظه‌کاری) و پین (پدر رادیکالیسم) دو قطب بحث شدند.

\textbf{۵. گسترش در قرن نوزدهم:} مفهوم چپ و راست در سراسر اروپا گسترش یافت و با ظهور سوسیالیسم، بُعد اقتصادی پیدا کرد.

\textbf{۶. تنوع درونی:} از همان آغاز، نه چپ یکدست بود و نه راست؛ اختلافات درونی همیشه وجود داشته‌اند.

\textbf{۷. محدودیت‌ها:} مدل خطی ساده‌سازی می‌کند، اما همچنان پرکاربرد است.

\end{summarybox}

\vspace{20pt}

\begin{questionbox}

\begin{enumerate}[nosep, rightmargin=1cm]
    \item آیا چینش تصادفی مجلس ۱۷۸۹ می‌توانست متفاوت باشد؟ اگر هواداران شاه در چپ می‌نشستند، تاریخ زبان سیاسی چگونه می‌شد؟

    \item برک و پین هر دو مدافع آزادی بودند—یکی از انقلاب آمریکا دفاع کرد (برک) و دیگری نقش مهمی در آن داشت (پین). چرا درباره انقلاب فرانسه اینقدر متفاوت فکر می‌کردند؟

    \item آیا تجربه دوران ترور، نقدی منصفانه به ایده‌های چپ است؟ یا هر انقلابی—چپ یا راست—ممکن است به خشونت بیانجامد؟

    \item با توجه به تنوع درونی چپ و راست، آیا اصلاً استفاده از این واژگان معنادار است؟
\end{enumerate}

\end{questionbox}

% ══════════════════════════════════════════════════════════════════════════════
%                         پایان فصل اول
% ══════════════════════════════════════════════════════════════════════════════